\begin{fmexfill}{Ejemplos: leemos los límites en la gráfica}
\begin{minipage}[t][\tcbtextheight][s]{\linewidth}
\noindent \begin{minipage}[t]{0.48\linewidth}\centering $f(x)=\dfrac{2x^2}{x^2+1}$\par\smallskip \begin{tikzpicture}\begin{axis}[fmaxis,width=6.4cm,height=4.6cm,xmin=-6,xmax=6,ymin=-1,ymax=3.2]\addplot[fmcurve,domain=-6:6,samples=160]{2*x^2/(x^2+1)};\draw[fmasy] ({rel axis cs:0,0}|-{axis cs:0,2}) -- ({rel axis cs:1,0}|-{axis cs:0,2});\end{axis}\end{tikzpicture}\par\smallskip $\displaystyle\lim_{x\to+\infty}f(x)=2,\qquad\lim_{x\to-\infty}f(x)=2$\\ \vspace*{0.2cm} Asíntota horizontal $y=2$ por los dos lados.\end{minipage}\hfill\begin{minipage}[t]{0.48\linewidth}\centering $g(x)=2^x-1$\par\smallskip \begin{tikzpicture}\begin{axis}[fmaxis,width=6.4cm,height=4.6cm,xmin=-6,xmax=4,ymin=-2,ymax=6]\addplot[fmcurve,domain=-6:3,samples=160]{2^x-1};\draw[fmasy] ({rel axis cs:0,0}|-{axis cs:0,-1}) -- ({rel axis cs:1,0}|-{axis cs:0,-1});\end{axis}\end{tikzpicture}\par\smallskip $\displaystyle\lim_{x\to-\infty}g(x)=-1,\qquad\lim_{x\to+\infty}g(x)=+\infty$\\ \vspace*{0.2cm} Asíntota horizontal $y=-1$ solo por la izquierda.\end{minipage}\par
\vfill
\noindent \begin{minipage}[t]{0.48\linewidth}\centering $h(x)=\dfrac{x^3}{10}-x$\par\smallskip \begin{tikzpicture}\begin{axis}[fmaxis,width=6.4cm,height=4.6cm,xmin=-6,xmax=6,ymin=-6,ymax=6]\addplot[fmcurve,domain=-6:6,samples=160]{x^3/10-x};\end{axis}\end{tikzpicture}\par\smallskip $\displaystyle\lim_{x\to-\infty}h(x)=-\infty,\qquad\lim_{x\to+\infty}h(x)=+\infty$\\ \vspace*{0.2cm} No tiene asíntota horizontal.\end{minipage}\hfill\begin{minipage}[t]{0.48\linewidth}\centering $k(x)=\sin x$\par\smallskip \begin{tikzpicture}\begin{axis}[fmaxis,width=6.4cm,height=4.6cm,xmin=-10,xmax=10,ymin=-2,ymax=2]\addplot[fmcurve,domain=-10:10,samples=300]{sin(deg(x))};\end{axis}\end{tikzpicture}\par\smallskip $\displaystyle\nexists\lim_{x\to-\infty}\sin x,\qquad\nexists\lim_{x\to+\infty}\sin x$\\ \vspace*{0.2cm} Oscila entre $-1$ y $1$ sin parar.\end{minipage}\par
\vfill
\noindent \begin{minipage}[t]{0.48\linewidth}\centering $m(x)=\sqrt{x-2}$\par\smallskip \begin{tikzpicture}\begin{axis}[fmaxis,width=6.4cm,height=4.6cm,xmin=-4,xmax=10,ymin=-1,ymax=4]\addplot[fmcurve,domain=2:10,samples=300]{sqrt(max(x-2,0))};\node[fmdot] at (axis cs:2,0) {};\end{axis}\end{tikzpicture}\par\smallskip $\displaystyle\nexists\lim_{x\to-\infty}m(x),\qquad\lim_{x\to+\infty}m(x)=+\infty$\\ \vspace*{0.2cm} Dominio $[2,+\infty)$: no está definida para $x<2$.\end{minipage}\hfill\begin{minipage}[t]{0.48\linewidth}\centering $n(x)=\sqrt{9-x^2}$\par\smallskip \begin{tikzpicture}\begin{axis}[fmaxis,width=6.4cm,height=4.6cm,xmin=-5,xmax=5,ymin=-1,ymax=4]\addplot[fmcurve,domain=-3:3,samples=300]{sqrt(max(9-x^2,0))};\node[fmdot] at (axis cs:-3,0) {};\node[fmdot] at (axis cs:3,0) {};\end{axis}\end{tikzpicture}\par\smallskip $\displaystyle\nexists\lim_{x\to-\infty}n(x),\qquad\nexists\lim_{x\to+\infty}n(x)$\\ \vspace*{0.2cm} Dominio $[-3,3]$: no llega a ninguno de los dos infinitos.\end{minipage}\par
\vfill

      
\begin{fmnote}\fmuline{Truco}: para \fmm{x\to+\infty} mira \textbf{el final derecho} de la gráfica; para \fmm{x\to-\infty}, \textbf{el final izquierdo}. Una tabla de valores con \fmm{x=10,\ 100,\ 1000,\dots} también te lo dice.\end{fmnote}
\end{minipage}
\end{fmexfill}
